\documentclass{article}
\usepackage{amsmath,amssymb,xurl}
\usepackage[hidelinks]{hyperref}
% Shared commands for notes in docs/paper-gaps/.

\DeclareUnicodeCharacter{2080}{\ifmmode{}_{0}\else\textsubscript{0}\fi}
\DeclareUnicodeCharacter{2081}{\ifmmode{}_{1}\else\textsubscript{1}\fi}
\DeclareUnicodeCharacter{2082}{\ifmmode{}_{2}\else\textsubscript{2}\fi}
\DeclareUnicodeCharacter{2099}{\ifmmode{}_{n}\else\textsubscript{n}\fi}

\newcommand{\Lean}{\textsc{Lean}}
\ExplSyntaxOn
\NewDocumentCommand{\leanid}{m}{
  \ifmmode
    \textsf{\detokenize{#1}}
  \else
    \group_begin:
    \urlstyle{sf}
    \tl_set:Nn \l_tmpa_tl {#1}
    \tl_replace_all:Nnn \l_tmpa_tl {\_} {_}
    \exp_args:NV \path \l_tmpa_tl
    \group_end:
  \fi
}
\ExplSyntaxOff
\providecommand{\tr}{\operatorname{tr}}
\newcommand{\tnleanIssueBase}{https://github.com/LionSR/TNLean/issues/}
\newcommand{\tnleanPrBase}{https://github.com/LionSR/TNLean/pull/}
\newcommand{\tnleanDocBase}{https://sirui-lu.com/TNLean/docs/}
\newcommand{\ghissue}[1]{\href{\tnleanIssueBase#1}{GitHub issue \##1}}
\newcommand{\ghpr}[1]{\href{\tnleanPrBase#1}{PR \##1}}
\newcommand{\doclink}[1]{\href{\tnleanDocBase#1.html}{\path{#1}}}

% Dirac notation and the identity operator, as in blueprint/src/macros/common.tex.
% \providecommand keeps note-local definitions authoritative.
\usepackage{dsfont}
\providecommand{\ket}[1]{|#1\rangle}
\providecommand{\bra}[1]{\langle #1|}
\providecommand{\braket}[2]{\langle #1 | #2 \rangle}
\providecommand{\ketbra}[2]{|#1\rangle\!\langle #2|}
\providecommand{\Id}{\mathds{1}}
\providecommand{\one}{\mathds{1}}
\providecommand{\C}{\mathbb{C}}
\providecommand{\MN}[1]{M_{#1}(\C)}
\providecommand{\MD}{M_D(\C)}

% Verdict marker: \gapnote{<kind>}{<status>}. Kind is the policy
% classification (clarification, local-correction, scope-restriction,
% unfaithful, false-source, open-gap); status is open, wip (elimination
% underway), resolved, or historical. Machine-read by the site index;
% prints a verdict line.
\newcommand{\gapnote}[2]{\par\noindent\textbf{Verdict:} #1 (#2).\par}

\title{A Physical Wilson-Loop Obstruction to Equal Local Density}
\author{}
\date{}
\begin{document}
\maketitle
\gapnote{open-gap}{open}

Schuch, Cirac, and P\'erez-Garc\'ia,
\cite[Lemma~6.14 and Theorem~6.15]{Schuch2010PEPS},
source \path{Papers/1001.3807/paper_v3.tex}, lines 2199--2267,
describe flux-string deformation and detection of a single fluxon's type.
Theorem 6.15 measures a conjugate $xgx^{-1}$ on the supporting plaquette,
so it determines the conjugacy class $[g]$. The assertion below is a
consequence of regular isometry for general observed closed walks; it is
not a literal special case of the source's particular plaquette measurement.

Let $\Gamma$ be a finite graph, $G$ a finite group and $a$ a family of regular
$G$-isometric physical sites. Let $R$ be an observation region whose induced
graph is connected, and let $\gamma:o\to o$ be a finite closed walk entirely
in that induced graph. For an assignment $u$ of oriented bond insertions,
let $H_u(\gamma)$ be its ordered transport around $\gamma$, with inverse
transport on reversed edges. Write $[H_u(\gamma)]$ for its conjugacy class.
Let $M_u(\sigma,\tau)$ be the actual physical coefficient matrix across
$R$ and $R^c$, with the insertion $u$ on the original bond contraction.
For two assignments $u,v$, the proved assertion is
\begin{align}
    [H_u(\gamma)]\ne[H_v(\gamma)]
    \ \Longrightarrow\
    M_u\ne0,\quad M_v\ne0,\quad
    \frac{M_uM_u^\dagger}{\operatorname{tr}(M_uM_u^\dagger)}
    \ne
    \frac{M_vM_v^\dagger}{\operatorname{tr}(M_vM_v^\dagger)}.
    \label{eq:scp10_wilson_density_separation}
\end{align}
There is no connectedness assumption on $R^c$ and no simple-connectedness
assumption on $R$. The walk and the two computed holonomy classes are explicit
inputs; no physical measurement or density inequality is supplied.
\footnote{\leanid{TNLean.PEPS.regularPhysicalCutMatrix_normalized_reduced_ne_of_closedWalk_holonomy}.}

The complete closed-walk measurement is constructed from local regular
isometry. Its projector $Q$ onto the class $[H_u(\gamma)]$ acts on every
open boundary column, hence on the actual cuts:
\begin{align}
    QM_u=M_u,\qquad QM_v=0.
    \label{eq:scp10_wilson_density_projector}
\end{align}
Local physical inverses prove both cuts nonzero. Put
$\rho_u=M_uM_u^\dagger/\operatorname{tr}(M_uM_u^\dagger)$ and define
$\rho_v$ similarly. Equation~(\ref{eq:scp10_wilson_density_projector}) gives
$Q\rho_u=\rho_u$ and $Q\rho_v=0$. Thus $\rho_u=\rho_v$ would force
$\rho_u=0$, contradicting $\operatorname{tr}\rho_u=1$.

On a native torus with both periods at least three, the native bond labels
are converted to directed graph transport and the native physical cut
equals the graph cut above. The same implication therefore holds for
actual native bond insertions.
\footnote{\leanid{TNLean.PEPS.IsGIsometric.torusBond_normalized_density_ne_of_closedWalk_holonomy}.}
Taking $v=1$ supplies a sufficient criterion for failure of vacuum density.

This is a sufficient obstruction, not a classification of annular
observations or a converse density theorem. Deriving a suitable observed
walk with nontrivial holonomy from a proposed endpoint placement remains
a geometric obligation. In particular, endpoint exclusion alone supplies
no vanishing-holonomy premise. This physical obstruction is distinct from
the nonabelian seam-word obstruction to restricted homotopy, whose
coefficient group need not be the physical insertion group.
The broader source interpretation remains open under \ghissue{8703};
no Hamiltonian common-kernel result from \ghissue{8947} is used.

\bibliographystyle{alpha}
\bibliography{references}
\end{document}
